% ECONOMICS_PROBLEM: {"id":"C73-25","title":"Minimizer selection by the uncorrected vanishing-common-noise game","jel_code":"C73","source_id":"OP-0215"}
% New manuscript; earlier statements and research attempts are preserved.
% !TEX program = pdflatex
\documentclass[11pt,a4paper]{article}
\usepackage[T1]{fontenc}
\usepackage{lmodern}
\usepackage[left=20mm,right=20mm,top=18mm,bottom=18mm]{geometry}
\usepackage{amsmath,amssymb,amsthm,mathtools,microtype}
\usepackage[colorlinks=true,linkcolor=blue,citecolor=blue,urlcolor=blue,hypertexnames=false]{hyperref}
\allowdisplaybreaks
\setlength{\parindent}{0pt}
\setlength{\parskip}{3pt}
\setlength{\emergencystretch}{2em}
\newtheorem{theorem}{Theorem}
\newtheorem{lemma}[theorem]{Lemma}
\newcommand{\R}{\mathbb R}
\newcommand{\E}{\mathbb E}
\newcommand{\PP}{\mathbb P}
\newcommand{\F}{\mathcal F}
\newcommand{\G}{\mathcal G}
\newcommand{\Hb}{\overline H}
\newcommand{\tr}{\operatorname{tr}}
\newcommand{\Sym}{\operatorname{Sym}}
\newcommand{\TV}{\operatorname{TV}}
\newcommand{\one}{\mathbf 1}
\newcommand{\dt}{\,dt}
\title{Proof: A nonminimizing limit of the uncorrected noisy game}
\author{JiHo Bae}
\date{11 October 2026}
\begin{document}
\maketitle
\begin{abstract}
The answer to C73-25, with its stated admissible class, is negative.
We construct fixed smooth potential costs in three states and smooth,
nonnegative repulsions tending locally uniformly to zero. Their exact
uncorrected master equations have unique classical solutions. The actual
selected population laws converge to a deterministic path which is not
the population path of any global potential minimizer. The selection
argument is an area identity along the actual noisy population; it does
not identify a local characteristic graph with the global master field.
Boundary scales, tightness and removal of the repulsion are included.
\end{abstract}

\textbf{Scope.} The canonical statement permits a vanishing interior
perturbation of the boundary repulsion, with no monotonicity or support
restriction. The example uses that permission. It does not settle the
more restricted question in which the penalty must be a decreasing
cutoff supported only near zero. The original costs remain unchanged.
The cost-corrected selection theorem in
\cite[Theorem 4.5, p.~3692]{CD} and \cite{CeD} has different
equations and is not invoked as an uncorrected selection theorem.

% FIRST_POSTED: 2026-10-11
% ATTEMPT_STATUS: solved
% ENTRY_KIND: research_attempt
The complete original proof is preserved at the immutable
\href{https://github.com/jbaelaw/uncorrected-noise-nonminimizing-selection/blob/df1d03253a80b693cc7880155b0a3cf0a4c809e7/paper/Proof.tex}{Proof TeX}
and \href{https://github.com/jbaelaw/uncorrected-noise-nonminimizing-selection/blob/df1d03253a80b693cc7880155b0a3cf0a4c809e7/paper/Proof.pdf}{Proof PDF}.

\section{The equation and the conclusion}
Write $\Delta_3=\{m\geq0:\sum_i m_i=1\}$ and
$A_m=\operatorname{diag}(m)-mm^\top$. Set $T=1$ and $\F=0$.
For $U\in\R^3$ put
\[
 a_{ij}=(U_i-U_j)_+,
 \quad H_i(U)=\tfrac12\sum_{j\ne i}a_{ij}^2,
 \quad b_i^a=\sum_{j\ne i}(m_ja_{ji}-m_ia_{ij}).
\]
Given a penalty $\phi$, set
\[
 r_i^\phi=\phi(m_i)-m_i\sum_j\phi(m_j),\qquad
 Q_i=\sum_j\phi(m_j)(U_j-U_i),\qquad b=b^a+r^\phi.
\]
The exact master equation used below is
\begin{equation}\label{master}
 \partial_tU_i+\varepsilon A_m:D^2U_i
 +[b+2\varepsilon(e_i-m)]\cdot DU_i-H_i(U)+Q_i=0,
 \qquad U_i(1,m)=\partial_i\G(m).
\end{equation}
All derivatives are tangent derivatives. The selected population has
drift $b(m,U(t,m))$ and covariance $2\varepsilon A_m\dt$, with the
Brownian realization specified in C73-25. In particular, its noise is
the original Kimura noise.

\begin{theorem}\label{main}
There are a fixed smooth ambient potential $\G$, a fixed
$m_0\in\operatorname{int}\Delta_3$, and an admissible penalty family
$\phi_\varepsilon$ for which the actual solutions of \eqref{master}
and their selected populations satisfy
\[
 m^\varepsilon\longrightarrow \bar m
 \quad\hbox{in probability in }C([0,1],\Delta_3),
 \qquad \bar m\notin\mathcal M_{1,m_0}.
\]
Thus their limiting law $Q=\delta_{\bar m}$ satisfies
$Q(\mathcal M_{1,m_0})=0$.

More precisely, for one sufficiently small fixed $\delta>0$, the
construction has
\begin{equation}\label{scales}
 \eta_\varepsilon=\frac1{\log(e+1/\varepsilon)},\qquad
 \theta_\varepsilon=\min\{\varepsilon^2,1/100,1/(16\sqrt3)\},
 \qquad
 \phi_\varepsilon(s)=\kappa_\varepsilon
       \chi(s/\theta_\varepsilon)^2+\eta_\varepsilon\psi_\delta(s).
\end{equation}
Here $\chi$ is smooth, nonincreasing, one on $[0,1]$ and zero on
$[2,\infty)$; $\psi_\delta\geq0$ is fixed and smooth. A sufficient
height $\kappa_\varepsilon$ is specified in Lemma~\ref{admissible}.
There is no restriction on the product
$\kappa_\varepsilon\theta_\varepsilon$. In addition,
\[
 \E\int_0^1\|r^{\phi_\varepsilon}(m_t^\varepsilon)\|_1\dt
       \longrightarrow0,
\]
and the population laws are tight in the required path space.
\end{theorem}

Constants denoted $C$ are independent of the small parameters until
those parameters are fixed. A subscript $\delta$ permits dependence
on the fixed construction, including its small radial scale.
The order of choices is: fixed cutoffs, a small fixed $\rho_0$, the
absorption constant $K_0$ determined by the uniform profile and forward
Poincar\'e bounds, then a sufficiently small fixed $\delta$ and
$R=\delta^2$, and finally
$\eta=\eta_\varepsilon\to0$ and $\varepsilon\to0$.

\section{Exact flat coordinates in one ordering chamber}
On the chamber $U_1>U_2>U_3$ introduce
\begin{equation}\label{flat}
 y_1=\sqrt{2m_1},\qquad y_2=\sqrt{2m_1+4m_2},\qquad
 P_1=y_1\left(U_1-\tfrac{U_2+U_3}{2}\right),\qquad
 P_2=\tfrac{y_2}{2}(U_2-U_3).
\end{equation}
The inverse map is
\[
 m_1=y_1^2/2,\quad m_2=(y_2^2-y_1^2)/4,\quad
 m_3=1-(y_1^2+y_2^2)/4.
\]
Consequently $U\cdot dm=P\cdot dy$ and
\begin{equation}\label{flatH}
 \Hb(m,U):=\sum_i m_iH_i(U)=\tfrac12|P|^2,
 \qquad U_1>U_2>U_3\quad\Longleftrightarrow\quad P_1/y_1>P_2/y_2>0.
\end{equation}
Indeed, the three positive differences are
$P_1/y_1-P_2/y_2$, $P_1/y_1+P_2/y_2$, and $2P_2/y_2$.
Inserting them into $\sum_i m_iH_i$ proves \eqref{flatH}.
Thus the original deterministic feedback dynamics with zero running
cost, in this chamber and without repulsion, are
$\dot y=-P$, $\dot P=0$.

For an interior penalty $\phi=\eta\psi$, put
$\Psi=\sum_i\psi(m_i)$ and define
\begin{equation}\label{RD}
 \begin{split}
 R_1(y)&=\frac{\psi(m_1)}{y_1}-\frac\Psi2y_1,\\
 R_2(y)&=\frac{2\psi(m_2)+\psi(m_1)}{y_2}-\frac\Psi2y_2,\\
 D(y)&=\operatorname{diag}(\Psi+R_1/y_1,\Psi+R_2/y_2).
 \end{split}
\end{equation}
The exact deterministic equations of the uncorrected game are
\begin{equation}\label{ODE}
 \dot y=-P+\eta R(y),\qquad \dot P=\eta D(y)P.
\end{equation}
For example, with $J_1=y_1$, $J_2=y_2/2$ and
$p_{v,k}=P_k/J_k$, the equations for component differences give
$\dot p_{v,k}=c_kp_{v,k}^2/2+\eta\Psi p_{v,k}$,
where $c_1=2,c_2=1$. Also
$\dot J_k=-c_kP_k/2+\eta(R_k/y_k)J_k$.
The quadratic terms cancel in $\dot P_k$ and give \eqref{ODE}.
There is no derivative of $\psi$ in this costate equation.

\section{A fixed nonnegative profile with centered rotation}
Let
\[
 m_*=(1/5,3/10,1/2),\quad u=\sqrt{2/5},\quad y_*=(u,2u),
 \quad p=(1,1),\quad P_*=\delta p,\quad y_0=y_*+\delta p.
\]
Set $m_0=m(y_0)$, using the inverse map in Section~2.
Take disjoint fixed smooth cutoffs $\zeta_1,\zeta_2$, equal to one
near $1/5,3/10$ respectively and supported in small interior intervals
disjoint from $1/2$. For parameters $h_1,h_2,d_1,d_2$ set
\[
 h_a(s)=\sum_{j=1}^2[h_j+d_j(s-s_j)]\zeta_j(s),
 \qquad s_1=1/5,\quad s_2=3/10,\qquad
 \psi_{\delta,a}(s)=s+\delta h_a(s).
\]
The baseline $\psi(s)=s$ has $R=0$ and $D=I$ in \eqref{RD}.
Linearity therefore gives
\begin{equation}\label{baseline}
 R_{\delta,a}=\delta R_h,\qquad D_{\delta,a}=I+\delta D_h,
 \qquad M_h=D_h-DR_h.
\end{equation}
Here $DR_h$ means the ordinary Jacobian with respect to $y$.

Fix the two slopes at
\[
 d_1=-\frac5{7u},\qquad d_2=-\frac1{7u}.
\]
Choose the two values by the following exact linear equations:
\begin{equation}\label{Z}
 Z_\delta(a):=-\int_0^1
 \left[R_{\delta,a}(y_*+s\delta p)
 +(1-s)D_{\delta,a}(y_*+s\delta p)\delta p\right]ds=0.
\end{equation}
The map $Z_\delta/\delta$ is smooth at $\delta=0$, where it equals
$-R_h(y_*)-p/2$. Direct substitution in \eqref{RD} gives
\[
 \frac{R_{h1}(y_*)}{u}=2h_1-\tfrac12h_2,
 \qquad \frac{R_{h2}(y_*)}{2u}=\tfrac18h_1+\tfrac34h_2.
\]
The values-to-drift matrix has determinant $25u^2/8\ne0$.
Thus, for all sufficiently small $\delta$, inversion of a smooth
two-by-two matrix gives unique smooth values $h_1(\delta),h_2(\delta)$
satisfying \eqref{Z}, with limits
\begin{equation}\label{jets}
 h_1(0)=-\frac8{25u},\qquad h_2(0)=-\frac7{25u}.
\end{equation}
The functions $h_{a(\delta)}$ have uniform smooth bounds.
Outside their fixed interior supports $\psi_\delta(s)=s$;
on the supports $s$ is bounded below. Hence
$\psi_\delta=s+\delta h_{a(\delta)}\geq0$ on $[0,1]$ after reducing
$\delta$. This fixes one smooth scalar profile, rather than unrelated
jets prescribed at different times.

We record the matrix calculation responsible for the selection.
At $y_*$ let $K_h=15h_1/4-5h_2/2$. Differentiating \eqref{RD} yields
\[
 \begin{split}
 (M_h)_{11}-(M_h)_{22}&=K_h-4d_1/5+d_2/2,\\
 (M_h)_{12}+(M_h)_{21}&=-d_1/10+d_2/2,\\
 (M_h)_{12}&=d_2/5.
 \end{split}
\]
At the limiting values \eqref{jets} these give
\begin{equation}\label{Mmatrix}
 M_h(y_*)=\frac1u
 \begin{pmatrix}-113/70&-1/35\\1/35&-113/70\end{pmatrix}
 =\mu_0I+\omega_0J,
 \qquad J=\begin{pmatrix}0&1\\-1&0\end{pmatrix},\quad
 \omega_0=-\frac1{35u}<0.
\end{equation}
For our values at positive $\delta$, the error in this matrix is
$O(\delta)$. Freeze at $y_*$ the matrices
\begin{equation}\label{ABC}
 A=DR_\delta(y_*),\quad D=D_\delta(y_*),\quad B=D-A,
 \quad C=J\left(B-\tfrac12\tr(B)I\right),\quad
 k=\tr D-\tfrac12\tr B.
\end{equation}
Every traceless two-by-two matrix becomes symmetric after left
multiplication by $J$. By \eqref{baseline}--\eqref{Mmatrix},
\begin{equation}\label{coercivematrices}
 A=O(\delta),\quad D=I+O(\delta),\quad k=1+O(\delta),\quad
 c_0\delta I\leq C\leq c_1\delta I,\quad
 A^\top+D-kI=O(\delta)
\end{equation}
for positive fixed constants $c_0,c_1$. In fact
$C=\delta|\omega_0|I+O(\delta^2)$.

\section{Fixed smooth costs and a deterministic reference}
Let
\[
 c=(3\delta/(2u),\delta/u,0),\qquad
 \G_{\rm b}(m)=c\cdot m.
\]
Its flat gradient at $y_*$ is $P_*$. Choose once and for all a smooth
bounded function $q$ on $[0,\infty)$, equal to $-2$ on $[0,1]$ and
zero on $[4,\infty)$. With $R=\delta^2$ set
\begin{equation}\label{radial}
 H_R(r)=\int_0^r t q(t^2/R^2)\,dt.
\end{equation}
Then $H_R(r)=-r^2$ for $r\leq R$, its gradient vanishes for
$r\geq2R$, and
\begin{equation}\label{radialbounds}
 |\nabla H_R|\leq CR,\qquad
 \|D^2H_R\|+|H_R'(r)/r|\leq C,\qquad
 \|D^3H_R\|\leq C/R.
\end{equation}
The value at $r=0$ in the quotient is interpreted by continuity.
The composition is a smooth radial function of $y$.

On a flat ball of fixed radius $\rho_0$ about $y_*$ define the desired
local terminal potential
\begin{equation}\label{localG}
 \widetilde\G(y)=\G_{\rm b}(m_*)+P_*\cdot(y-y_*)+H_R(|y-y_*|).
\end{equation}
Multiply $\widetilde\G-\G_{\rm b}(m(y))$ by a smooth cutoff equal to
one on radius $\rho_0/2$ and zero outside radius $\rho_0$, and add
$\G_{\rm b}$. Choose the ball compactly inside the simplex chart.
This defines a smooth potential on the whole simplex. Its degree-one
homogeneous extension
$\G(m)=S\G(m/S)$, $S=\sum_i m_i>0$, defines a smooth ambient
potential on a neighborhood of the simplex. Since the linear flat
term and $\G_{\rm b}$ agree to first order at $y_*$, this extension has
\begin{equation}\label{gsmall}
 \max_i\|g_i-c_i\|_\infty\leq C(\delta\rho_0+R),\qquad
 K:=\max_i\|g_i\|_\infty\leq C\delta.
\end{equation}
All these costs are fixed before $\varepsilon$ varies.

For small $\eta\geq0$, solve \eqref{ODE} backward with terminal
position $z$ near $y_*$ and momentum $\nabla\widetilde\G(z)$.
At $\eta=0,z=y_*$ this is
$\bar y^0(t)=y_*+(1-t)P_*$, $\bar P^0(t)=P_*$.
The terminal-to-initial position derivative is
$I+D^2\widetilde\G(y_*)=-I$, hence is invertible.
Differentiating the backward ODE with respect to $\eta$ at this
reference gives exactly the initial-position derivative $Z_\delta$
in \eqref{Z}: in backward time $s$ the momentum derivative is
$-\int_0^sD_\delta(y_*+r\delta p)\delta p\,dr$, and its position
derivative is its integral minus $\int_0^sR_\delta(y_*+r\delta p)\,dr$.
Therefore \eqref{Z} and the implicit function theorem give an exact
deterministic reference $(\bar y^\eta,\bar P^\eta)$ with
\begin{equation}\label{reference}
 \bar y^\eta(0)=y_0,\qquad
 \bar y^\eta(1)=y_*+s_\eta,\quad |s_\eta|\leq C_\delta\eta^2,
 \qquad |\bar y^\eta-y_*|+|\bar P^\eta|\leq C\delta.
\end{equation}
It converges uniformly to $(\bar y^0,P_*)$.
This reference is used for centering; no assertion about a global
characteristic graph is required.

\section{Admissibility and estimates independent of layer height}
\begin{lemma}\label{admissible}
Let $q_0\geq0$ be any fixed smooth interior addition with
$\|q_0\|_\infty\leq Q$. For fixed $0<\varepsilon<1/2$ and
$0<\theta<1/(4\sqrt3)$, the penalty
$\phi(s)=\kappa\chi(s/\theta)^2+q_0(s)$ is admissible for
\eqref{master} if
\begin{equation}\label{height}
 \kappa\geq\max\{\kappa_H(\sqrt{2\varepsilon},2K+Q+2\varepsilon),
                     1+\varepsilon\}.
\end{equation}
Here $\kappa_H(\sigma,B_0)$ is the sufficient threshold of
\cite[Theorem 2.10]{BCCD}. It depends on noise and the bound $B_0$
for the nonnegative inward coefficient, and not on $\theta$ or the
large outflow bound. The solution is unique in the canonical classical
class and satisfies $|U_i|\leq K$. The same statement holds on any
fixed finite horizon, with $K=\|g\|_\infty$ since $\F=0$.
\end{lemma}
\begin{proof}
We give the reduction to the cited linear theory, since the printed
master theorem \cite[Theorem 3.4]{BCCD} assumes a decreasing cutoff
and cannot directly be applied to our total penalty.
Put $a_M(r)=\min(M,r_+)$, $M=2K$, and use
$H_M^{-,i}(v)=-\sum_{j\ne i}\int_0^{(v_i-v_j)_+}\min(M,r)\,dr$.
With $\phi_0=\kappa\chi(s/\theta)^2$, the generic fixed-point system
of \cite[Section 4.3.2, (4.16)--(4.18)]{BCCD} has component drift
$\phi_0(m_j)+\beta_j^i+m_j\beta_j^\circ$, where we insert
\[
 \begin{split}
 \beta_j^i&=q_0(m_j)+\sum_\ell m_\ell a_M(v_\ell-v_j)
                         +2\varepsilon\delta_{ij},\\
 \beta_j^\circ&=-\sum_\ell[\phi(m_\ell)+a_M(v_j-v_\ell)]
                         -2\varepsilon.
 \end{split}
\]
These coefficients sum to zero in the complete drift and give precisely
the clipped population drift plus $2\varepsilon(e_i-m)$.
Crucially, $0\leq\beta_j^i\leq Q+M+2\varepsilon$; the large
$\kappa$ appears in the outflow coefficient only. If
$P_\Gamma$ clips to $[-\Gamma,\Gamma]$, use the bounded source
\[
 P_\Gamma\left[H_M^{-,i}(v)
                +\sum_j\phi(m_j)P_\Gamma(v_j-v_i)\right].
\]
For fixed parameters all the coefficient regularity assumptions of
that construction hold. Its frozen linear solvability
\cite[Lemma 4.4]{BCCD}, the H\"older estimate
\cite[Theorem 2.10]{BCCD}, and the Schauder argument in its
Section 4.3.2 yield a Wright--Fisher classical solution of this clipped
system, with bounded first derivatives. Their constants may depend on
the fixed outflow and source sizes; the repulsion threshold uses the
displayed inward bound. This is why \eqref{height} is noncircular.

At a maximum over all components, $H_M^{-,i}\leq0$ and the switching
source is nonpositive; at a minimum $H_M^{-,i}=0$ and the switching
source is nonnegative. The maximum principle is justified at the
degenerate boundary using
$B_{\rm ent}=\log3+\sum_jm_j\log m_j$.
For any positive $\alpha$, a maximum of $U_i-\alpha B_{\rm ent}$
cannot lie on the boundary: an inward displacement changes $U_i$ by
$O(h)$ and $-\alpha B_{\rm ent}$ by a positive multiple of
$-\alpha h\log h+O(h)$. The component generator applied to
$B_{\rm ent}$ is bounded above for fixed parameters, since
$\varepsilon A_m:D^2B_{\rm ent}=2\varepsilon$ and the other entropy
drifts have finite upper bounds. Subtract also this bound plus one
times $\alpha(1-t)$, and apply the interior maximum principle in
reverse time. Letting $\alpha\downarrow0$ gives $|U_i|\leq K$.
Thus $M=2K$ removes the control and Hamiltonian clips. After $\kappa$
is fixed choose
$\Gamma>\max\{2K,4K^2+6K\|\phi\|_\infty\}$; both source clips
are then inactive. This proves existence for the exact unchanged
costs. If $K=0$, $U=0$ works, or one can use any positive control clip.

For bounded feedback rates $a_{ij}\leq M_1$,
$b_j\geq\phi(m_j)-C_\varepsilon m_j$.
On $m_j\leq\theta$, the generator of the population, or of any
component diffusion with added drift $2\varepsilon(e_i-m)$, applied
to $-\log m_j$ is at most
$-(\kappa-\varepsilon)/m_j+C_\varepsilon+2\varepsilon$.
Off that layer it is bounded above. Stopped logarithmic barriers
exclude hitting any face. Interior local Lipschitz coefficients give
a unique strong solution extending throughout $[0,1]$.

Finally, if $U,V$ are canonical classical solutions, their difference
$w_i$ satisfies a linear equation with the component generator based
on $U$ and source
\[
 \sum_j\phi(m_j)(w_j-w_i)-H_i(U)+H_i(V)
                  +[b(m,U)-b(m,V)]\cdot DV_i.
\]
For fixed $\varepsilon$ this source is bounded by
$C_\varepsilon\max_j\|w_j\|_\infty$, including the drift mismatch
term because $DV$ is bounded. It\^o's formula along the interior
component diffusion, followed by removal of compact stopping times,
gives backward Gronwall with zero terminal data. Hence $w=0$ in
precisely the stated canonical class.
\end{proof}

Apply the lemma in \eqref{scales} with
$q_0=\eta_\varepsilon\psi_\delta$ and
$Q=\|\psi_\delta\|_\infty$, taking $\kappa_\varepsilon$ at least
the right side of \eqref{height} and $1/\varepsilon$.
The noise convention in \cite{BCCD} is generator
$\sigma^2 A_m/2$, so $\sigma=\sqrt{2\varepsilon}$ here.
The cutoff is supported in $[0,2\theta_\varepsilon]$; hence the family
is smooth, nonnegative and locally uniformly vanishing on $(0,1]$.
Only small $\varepsilon$ matters. If an index for every
$\varepsilon>0$ is wanted, put $a=\max(1,4\varepsilon)$ and apply
the lemma to $V(s,m)=U(s/a,m)/a$ with horizon $a$, noise
$\varepsilon/a$, terminal data $g/a$ and penalty $\phi/a$;
rescaling back supplies the remaining indices without changing $g$.

Two useful uniform estimates also explain why an arbitrarily high
layer causes no hidden impulse. With $M=2K$, put
$\mu=\min_i m_{0i}$, $C_*=2(M+1)$ and
$\ell=(\mu/2)e^{-C_*}$. Let
$\alpha_\varepsilon=\sup_{s\in[\ell,1]}\phi_\varepsilon(s)\to0$.
Stop the actual population on its first coordinate reaching $\ell$.
Before this time $b_i+C_*m_i\geq0$, for small $\varepsilon$.
The integrating-factor martingale has expected squared supremum
at most $2\varepsilon e^{2C_*}$, since
$d\langle M_i^m\rangle\leq\varepsilon\dt/2$.
A coordinate exit needs a martingale drop of at least $\mu/2$.
Thus
\begin{equation}\label{globalexit}
 \PP(\hbox{exit before }1)\leq24\varepsilon e^{2C_*}/\mu^2.
\end{equation}
On nonexit paths the integrated repulsion is bounded by
$3\alpha_\varepsilon$. The controlled drift is uniformly bounded
and the population martingale tends uniformly to zero in $L^2$.
These facts prove tightness and repulsion removal in probability.

For expected absolute removal, set $F_{\rm ent}=\sum_i m_i\log m_i$
and $k_\ell=\log(1/(3\ell))>0$. Its repulsion drift is
\[
 r^\phi\cdot DF_{\rm ent}
 \leq-k_\ell\sum_{i:m_i<\ell}\phi(m_i)+3\alpha_\varepsilon\log3.
\]
Its controlled drift is at most $2M\log3$, its diffusion drift is
$2\varepsilon$, and its martingale bracket is bounded by
$24\varepsilon/e^2\dt$. Stop first inside the simplex and then pass
to the limit; the negative term is integrable by Fatou and the other
terms have the displayed bounds. On the exit event, apply the
resulting entropy inequality from the exit time to $1$ conditionally
on that event. The entropy oscillation is at most $\log3$.
Together with \eqref{globalexit}, this gives
\begin{equation}\label{impulse}
 \E\int_0^1\sum_i\phi_\varepsilon(m_i^\varepsilon)\dt
       \leq3\alpha_\varepsilon+C_\ell\varepsilon\longrightarrow0.
\end{equation}
Since $\|r^\phi\|_1\leq2\sum_i\phi(m_i)$, this proves the claimed
expected removal without a height-times-width assumption.

\section{The actual feedback remains in the protected chamber}
We now reduce $\rho_0$ and then $\delta$ as necessary.
The value estimate and \eqref{gsmall} give $|U_i|\leq C\delta$
globally, so all actual rates are $O(\delta)$.
For a component $i$, use the auxiliary diffusion with generator
\[
 \varepsilon A_m:D^2+[b(m,U)+2\varepsilon(e_i-m)]\cdot D.
\]
Start in the flat ball of radius $\rho_0/2$ and stop on radius
$\rho_0$. On this protected set the boundary cutoff is absent,
$\phi=\eta\psi_\delta$, and the drift is $O(\delta+\eta+\varepsilon)$.
Choose $\delta$ small compared with $\rho_0$ and then
$\eta,\varepsilon\leq\delta$. A deterministic drift displacement is
less than $\rho_0/4$, so Doob's inequality bounds the exit probability
by $C\varepsilon/\rho_0^2$.
Along this auxiliary diffusion \eqref{master} gives
$dU_i=(H_i-Q_i)\dt+$ martingale. Before exit,
$|H_i|\leq C\delta^2$ and $|Q_i|\leq C\eta\delta$.
Using \eqref{gsmall} at the terminal time, and the global value bound
at a stopped exit, proves uniformly on the inner ball and in time
\begin{equation}\label{chamberestimate}
 |U_i(t,m)-c_i|
 \leq C(\delta\rho_0+R+\delta^2+\eta\delta
                         +\delta\varepsilon/\rho_0^2).
\end{equation}
The gaps of $c$ are positive fixed multiples of $\delta$.
Choose $\rho_0$, then $\delta$ with $R=\delta^2$, and finally
$\eta,\varepsilon$ so that the right side is less than one quarter
of the smallest gap. This proves $U_1>U_2>U_3$ for the actual
global master field throughout the inner ball.

Let $L$ be a fixed sufficiently large constant and stop the actual
selected flat population at
\[
 \tau=1\wedge\inf\{t:|y_t-y_0|\geq L\delta\}.
\]
Reduce $\delta$ so that this ball and the reference
\eqref{reference} are inside radius $\rho_0/4$ about $y_*$.
The actual flat population drift is $O(\delta)$ and its bracket is
$O(\varepsilon)$. Taking $L$ greater than twice the deterministic
drift bound gives
\begin{equation}\label{localexit}
 \PP(\tau<1)\leq C\varepsilon/\delta^2.
\end{equation}
Up to this stop the actual and reference positions are $O(\delta)$
from $y_*$, their momenta are $O(\delta)$, and the exact flat and
radial formulas hold. The estimates concern the actual global
solution, not a separate reference diffusion.

\section{Original weighted energy and exact finite-noise errors}
For clarity we retain all derivative terms in the original equation.
Along the actual selected population write
$z_i=(e_i-m)\cdot DU_i$ and
$\Gamma_i=DU_i^\top A_mDU_i$. For
$V=\sum_i m_iU_i^2$, the generator including time gives
\begin{equation}\label{energyidentity}
 (\partial_t+\varepsilon A_m:D^2+b\cdot D)V
 =2\varepsilon\sum_i m_i\Gamma_i
   -\sum_{i,j\ne i}m_iU_j a_{ij}^2
   +\mathcal R,
 \quad
 \mathcal R=\sum_j\phi(m_j)\sum_i m_i(U_j-U_i)^2\geq0.
\end{equation}
Indeed, the cross term
$4\varepsilon U_i\Gamma_{A_m}(m_i,U_i)
=4\varepsilon m_iU_iz_i$ cancels the weighted-player contribution
$-4\varepsilon m_iU_iz_i$. On an active edge,
$a_{ij}(U_j^2-U_i^2)+U_ia_{ij}^2=-U_ja_{ij}^2$.
The repulsion terms combine as
$\sum_i r_iU_i^2-2\sum_i m_iU_iQ_i=\mathcal R$.
These identities prove \eqref{energyidentity} without dropping
any stochastic derivative correction.

Stop on $\min_i m_i\geq1/n$, apply It\^o's formula and take expectations.
The right side after rearrangement is bounded by
$K^2+2K(2K)^2$. The nonnegative left integrands pass by monotone
convergence, and the bounded remaining terms pass by bounded
convergence, because each continuous interior path stays in some
compact interior set. Hence
\begin{equation}\label{energybound}
 2\varepsilon\E\int_0^1\sum_i m_i\Gamma_i\dt
          +\E\int_0^1\mathcal R\dt\leq C_E,
 \qquad C_E=K^2+8K^3=O(\delta^2).
\end{equation}
In particular, on the actual stopped protected set,
$\varepsilon\E\int_0^\tau|D_{\rm tan}U|^2\dt\leq C_\delta$ and
$\varepsilon\E\int_0^\tau|D_{\rm tan}U|\dt\leq C_\delta\sqrt\varepsilon$.

Here are the precise transformed equations. Let $B_v$ have rows
$(1,-1/2,-1/2)$ and $(0,1,-1)$, so that $p_v=B_vU$.
Let $M^m$ be the population martingale,
$K_y=Dy\,A_m(Dy)^\top$, and put
\[
 \ell_k=\sum_i(B_v)_{ki}z_i,\qquad
 t_k=\sum_i(B_v)_{ki}DU_i^\top A_mD_my_k.
\]
On $[0,\tau]$, the exact equations are
\begin{equation}\label{stochasticflat}
 \begin{split}
 dy_k&=[-P_k+\eta R_k-\varepsilon K_{y,kk}/y_k]\dt+dM_k^y,\\
 dP_k&=[\eta D_{kk}(y)P_k+\varepsilon e_k]\dt+dN_k^P,\\
 e_k&=-K_{y,kk}P_k/y_k^2-2J_k\ell_k+c_kt_k,\\
 dM_k^y&=D_my_k\cdot dM^m,\\
 dN_k^P&=J_k\sum_i(B_v)_{ki}DU_i\cdot dM^m
                    +(P_k/y_k)D_my_k\cdot dM^m.
 \end{split}
\end{equation}
To check them, the original equation along the population is
$dU_i=(H_i-Q_i-2\varepsilon z_i)\dt+DU_i\cdot dM^m$.
Thus
$dp_{v,k}=(c_kp_{v,k}^2/2+\eta\Psi p_{v,k}
-2\varepsilon\ell_k)\dt+\sum_i(B_v)_{ki}DU_i\cdot dM^m$.
In $d(J_kp_{v,k})$ the quadratic drift cancels as before, the
coordinate It\^o drift gives the first term of $e_k$, and the product
bracket gives $\varepsilon c_kt_k$. This proves
\eqref{stochasticflat}, including the weighted-player and Kimura terms.
Bounded geometry on the protected set and \eqref{energybound} imply
\begin{equation}\label{noisebounds}
 \begin{split}
 \TV(E_X)&\leq C_\delta\varepsilon,\qquad
 \E\TV(E_p)\leq C_\delta\sqrt\varepsilon,\\
 \E\tr\langle M^y\rangle_\tau&\leq C_\delta\varepsilon,\qquad
 \E\tr\langle N^P\rangle_\tau\leq C_\delta,
 \end{split}
\end{equation}
where $dE_{X,k}=-\varepsilon K_{y,kk}/y_k\dt$ and
$dE_{p,k}=\varepsilon e_k\dt$.
In particular the costate martingale is bounded in energy; it is not
assumed to vanish. Cauchy--Schwarz for covariations gives
\begin{equation}\label{crossnoise}
 \E\TV\langle M_i^y,N_j^P\rangle_\tau
       \leq C_\delta\sqrt\varepsilon.
\end{equation}

\section{An area estimate for the actual centered processes}
Center the actual variables at the reference with $\eta=\eta_\varepsilon$:
$X=y-\bar y^\eta$, $p_c=P-\bar P^\eta$, stopping both at the same
$\tau$. In this section abbreviate $p_c$ to $v$.
Equations \eqref{reference}, \eqref{baseline} and Taylor's theorem
give the exact decomposition
\begin{equation}\label{centered}
 \begin{split}
 dX&=[-v+\eta AX+\eta r_X]\dt+dE_X+dM^y,\\
 dv&=[\eta Dv+\eta r_v]\dt+dE_p+dN^P,\\
 |r_X|&\leq C\delta^2|X|,\qquad
 |r_v|\leq C\delta^2(|X|+|v|),\qquad X_0=0.
 \end{split}
\end{equation}
For the second remainder, explicitly write
$[D(y)-D(y_*)]v+[D(y)-D(\bar y)]\bar P$.
Both derivatives of the nonconstant part of $D$ are $O(\delta)$;
the positions and $\bar P$ are $O(\delta)$.
The first remainder follows similarly from
$DR(y)-DR(y_*)=O(\delta|y-y_*|)$.
This explains the essential $\delta^2$ sizes.

On the good event $G=\{\tau=1\}$, the terminal radial formula gives
\begin{equation}\label{terminal}
 v_1=h_1X_1+e_1,\qquad |h_1|\leq H,\qquad
 |e_1|\leq C_\delta\eta^2|X_1|.
\end{equation}
Indeed, $v_1=\nabla H_R(X_1+s_\eta)-\nabla H_R(s_\eta)$.
Take $h_1=H_R'(|X_1|)/|X_1|$. The mixed difference from
$\nabla H_R(X_1)$ is bounded by
$\|D^3H_R\|\,|s_\eta|\,|X_1|$.
Use \eqref{radialbounds} and \eqref{reference}; $H$ is uniform in
small $\delta$, whereas $C_\delta$ is allowed to depend on $R$.

First, the forward equation in \eqref{centered}, Gronwall, Doob and
Cauchy--Schwarz give
\begin{equation}\label{Poincare}
 \E\sup_{t\leq1}|X_{t\wedge\tau}|^2
       \leq C\E\int_0^\tau|v_t|^2\dt+C_\delta\varepsilon.
\end{equation}
For example, its square root follows from the pathwise bound
$\sup|X|\leq2(\int_0^\tau|v|\dt+\TV(E_X)+\sup|M^y|)$,
valid once $\eta\|A\|+C\eta\delta^2$ is small.
It also bounds $\E\int_0^\tau w|X|^2\dt$ by
$C\E\int_0^\tau w|v|^2\dt+C_\delta\varepsilon$ for
$w(t)=e^{-\eta kt}$, since $1/2\leq w\leq1$ for small $\eta$.

Choose $K_0$ large, put $\lambda=K_0\delta^2$ and set
\[
 N=J+\eta\lambda I,\quad Q=X^\top Nv,\quad
 E=\tfrac12X^\top CX,\quad L_c=A^\top+D-kI.
\]
The two-by-two identity $A^\top J+JD=C+kJ$, together with It\^o's
formula, gives
\[
 \begin{split}
 dQ-\eta kQ\dt={}&[\eta X^\top Cv-\eta\lambda|v|^2
                   +\eta^2\lambda X^\top L_cv]\dt\\
 &+\eta[r_X^\top Nv+X^\top Nr_v]\dt+d\mathcal V_Q+d\mathcal M_Q,\\
 dE={}&[-X^\top Cv+\eta X^\top CAX+\eta X^\top Cr_X]\dt
                           +d\mathcal V_E+d\mathcal M_E,
 \end{split}
\]
where the complete finite-variation noise terms are
\[
 \begin{split}
 d\mathcal V_Q&=(dE_X)^\top Nv+X^\top N\,dE_p
                  +\sum_{ij}N_{ij}\,d\langle M_i^y,N_j^P\rangle,\\
 d\mathcal V_E&=X^\top C\,dE_X
                  +\tfrac12\sum_{ij}C_{ij}\,d\langle M_i^y,M_j^y\rangle.
 \end{split}
\]
All martingale terms have mean zero after stopping, by bounded
$X,v$ and \eqref{noisebounds}. The covariation term has been retained.
Integrating with factor $w$, substituting the energy equation and using
$X_0=0$ gives the exact identity
\begin{equation}\label{areaidentity}
 \begin{split}
 &\E[w(\tau)(Q_\tau+\eta E_\tau)]
     +\eta\lambda\E\int_0^\tau w|v|^2\dt\\
 &\quad+\eta^2\E\int_0^\tau
       wX^\top[\tfrac{k}{2}C-\Sym(CA)]X\dt\\
 ={}&\eta^2\lambda\E\int_0^\tau wX^\top L_cv\dt
   +\eta\E\int_0^\tau w[r_X^\top Nv+X^\top Nr_v]\dt\\
 &\quad+\eta^2\E\int_0^\tau wX^\top Cr_X\dt
    +\E\int_0^\tau w\,d\mathcal V_Q
    +\eta\E\int_0^\tau w\,d\mathcal V_E.
 \end{split}
\end{equation}
By \eqref{noisebounds}--\eqref{crossnoise}, the last two terms in
absolute value are at most $C_\delta\sqrt\varepsilon$.
By \eqref{coercivematrices}, after reducing $\delta$,
\[
 \tfrac{k}{2}C-\Sym(CA)\geq c_2\delta I.
\]
The first-order remainder on the right of \eqref{areaidentity} is
bounded by $C\eta\delta^2\E\int w(|X|^2+|X||v|)\dt$.
Use \eqref{Poincare} and $2|X||v|\leq|X|^2+|v|^2$.
Choosing $K_0$ sufficiently large absorbs it into one quarter of
$\eta\lambda\E\int w|v|^2\dt$, with error $C_\delta\varepsilon$.
The $L_c=O(\delta)$ term is absorbed into another quarter by the same
bound and small $\eta$. The remaining deterministic term is at most
$C\eta^2\delta^3\E\int w|X|^2\dt$ and is absorbed by the positive
bulk matrix for small $\delta$.

On $G$, \eqref{terminal} yields
\[
 Q_1+\eta E_1
 \geq\eta\big[c_0\delta/2-\lambda H
                   -C_\delta\eta-C_\delta\lambda\eta^2\big]|X_1|^2
 \geq c_3\eta\delta|X_1|^2.
\]
Here choose $\delta$ so that $\lambda H\leq c_0\delta/8$, then
$\eta$ small for that fixed $\delta$. On exit, $|X|+|v|\leq C\delta$;
the negative endpoint contribution has expectation at most
$C\delta^2\PP(\tau<1)\leq C\varepsilon$ by \eqref{localexit}.
Consequently \eqref{areaidentity} proves
\begin{equation}\label{selectionestimate}
 \eta\delta\E[\one_G|X_1|^2]
       +\eta\delta^2\E\int_0^\tau|v|^2\dt
             \leq C_\delta\sqrt\varepsilon.
\end{equation}
This estimate concerns the actual selected occupation law.
Since $\sqrt\varepsilon/\eta_\varepsilon\to0$, \eqref{Poincare},
\eqref{localexit} and \eqref{selectionestimate} imply
$\sup_t|y_t-\bar y^{\eta_\varepsilon}_t|\to0$ in probability.
The reference convergence in \eqref{reference} proves the path
convergence in Theorem~\ref{main}.

\section{The limiting population path is not minimizing}
The limit is the original central straight path
$\bar y(t)=y_*+(1-t)P_*$, with $\bar m=m(\bar y)$.
Its feedback stays in the strict chamber: at $y_*$,
$P_{*1}/y_{*1}=\delta/u>\delta/(2u)=P_{*2}/y_{*2}>0$,
and the inequalities persist along this short path.

It is necessary to exclude every control producing this population
path, not just its equilibrium feedback. For any nonnegative rates
$a$ and cotangent $U$, completing the square edge by edge gives
\begin{equation}\label{Fenchel}
 L(m,a)+U\cdot b^a(m,a)+\Hb(m,U)\geq0,
 \qquad L(m,a)=\tfrac12\sum_{i,j\ne i}m_i a_{ij}^2.
\end{equation}
Indeed each edge contributes
$m_i[\tfrac12a_{ij}^2-a_{ij}(U_i-U_j)
+\tfrac12(U_i-U_j)_+^2]\geq0$.
Along the central path choose the cotangent corresponding to $P_*$
in \eqref{flat}. If any rates generate that same path, then
$U\cdot\dot m=P_*\cdot\dot{\bar y}=-|P_*|^2$ and
$\Hb=|P_*|^2/2$. Thus \eqref{Fenchel} gives
$L\geq|P_*|^2/2$ almost everywhere. The central feedback attains
this lower bound.

Take any fixed $w\in\R^2$ with $|w|=R/2$ and use the straight
competitor from $y_0$ to $y_*+w$ with constant momentum $P_*-w$.
For sufficiently small fixed $\delta$, $R=\delta^2$ ensures that
this entire path lies in the protected strict chamber, so its original
rates are nonnegative and its kinetic action is $|P_*-w|^2/2$.
Its terminal potential is, by \eqref{localG} and \eqref{radial},
$\G_{\rm b}(m_*)+P_*\cdot w-|w|^2$.
Its total original action minus the minimum action among controls
generating the central population path is therefore
\[
 \tfrac12|P_*-w|^2+P_*\cdot w-|w|^2
                    -\tfrac12|P_*|^2=-\tfrac12|w|^2<0.
\]
This is a fixed positive action gap, independent of
$\varepsilon$ and $\eta$. No control producing $\bar m$ can be a
global minimizer. Together with the actual path convergence, this
proves Theorem~\ref{main} for the complete canonical quantifiers.
\qed

\begin{thebibliography}{9}
\bibitem{BCCD} E. Bayraktar, A. Cecchin, A. Cohen and F. Delarue,
\emph{Finite state mean field games with Wright--Fisher common noise},
Journal de Math\'ematiques Pures et Appliqu\'ees 147 (2021), 98--162.
DOI: \href{https://doi.org/10.1016/j.matpur.2021.01.003}{10.1016/j.matpur.2021.01.003}.
Version inspected: arXiv:1912.06701v2 (10 January 2021).
Theorem 2.10, journal p.~112, Lemma 4.4, p.~134, and Section 4.3.2,
p.~136, especially (4.16)--(4.18), supply the analytic construction.
\url{https://arxiv.org/abs/1912.06701v2}.
\bibitem{CeD} A. Cecchin and F. Delarue,
\emph{Selection by vanishing common noise for potential finite state
mean field games}, Communications in Partial Differential Equations
47(1) (2022), 89--168.
DOI: \href{https://doi.org/10.1080/03605302.2021.1955256}{10.1080/03605302.2021.1955256}.
\url{https://arxiv.org/abs/2005.12153}.
\bibitem{CD} P. Cardaliaguet and F. Delarue,
\emph{Selected topics in mean field games}, Proc. ICM 2022, vol.~5,
pp.~3660--3703, EMS Press (2023).
DOI: \href{https://doi.org/10.4171/ICM2022/17}{10.4171/ICM2022/17}.
The uncorrected question is in Section 4.3, p.~3691; (4.14), p.~3687,
is the master equation, and Theorem 4.5 is the cost-corrected result.
\url{https://ems.press/content/book-chapter-files/33265}.
\end{thebibliography}
\end{document}
